% =============================================================================
% SYSTEMATIC INTERVENTION DESIGN: A BEHAVIORAL ECONOMICS FRAMEWORK
% =============================================================================
% Style: Ernst Fehr (AER/JPE Academic Format)
% Based on: EBF Chapter 17 (Intervention Foundations)
% =============================================================================

\documentclass[12pt,a4paper]{article}

% -----------------------------------------------------------------------------
% PACKAGES
% -----------------------------------------------------------------------------
\usepackage[utf8]{inputenc}
\usepackage[T1]{fontenc}
\usepackage[english]{babel}
\usepackage{geometry}
\usepackage{setspace}
\usepackage{amsmath,amssymb,amsthm}
\usepackage{graphicx}
\usepackage{booktabs}
\usepackage{array}
\usepackage{natbib}
\usepackage{hyperref}
\usepackage{enumitem}
\usepackage{float}

% -----------------------------------------------------------------------------
% PAGE SETUP (AER Style)
% -----------------------------------------------------------------------------
\geometry{margin=2.5cm, top=3cm, bottom=3cm}
\setstretch{1.5}
\setlength{\parindent}{1.5em}
\setlength{\parskip}{0pt}

% Theorem environments
\newtheorem{proposition}{Proposition}
\newtheorem{axiom}{Axiom}
\newtheorem{definition}{Definition}
\newtheorem{corollary}{Corollary}

% -----------------------------------------------------------------------------
% TITLE
% -----------------------------------------------------------------------------
\title{\Large\textbf{Systematic Intervention Design:\\A Behavioral Economics Framework}}

\author{
Evidence-Based Framework Project\thanks{This paper synthesizes the intervention design methodology from the Evidence-Based Framework for Economic and Social Behavior (EBF). Corresponding documentation: Appendix HHH (METHOD-TOOLKIT). We thank Ernst Fehr, Urs Fischbacher, and colleagues at the University of Zurich for foundational contributions to behavioral economics methodology.}
}

\date{January 2026}

% =============================================================================
% DOCUMENT
% =============================================================================
\begin{document}

\maketitle

% -----------------------------------------------------------------------------
% ABSTRACT
% -----------------------------------------------------------------------------
\begin{abstract}
\noindent Behavioral interventions frequently underperform expectations, with failure rates exceeding 60\% in applied settings. We argue that this systematic underperformance results from ad-hoc intervention design that neglects context dependency, journey phase alignment, and complementarity structures. This paper presents a formal framework for systematic intervention specification based on three core contributions: (i) a taxonomy of nine intervention types distinguished by their underlying behavioral mechanisms, (ii) a 20-field schema for standardized intervention documentation enabling organizational learning, and (iii) five axioms governing effective intervention design. We demonstrate that intervention effectiveness depends critically on the match between intervention type and the target population's position in the behavioral change journey. Furthermore, we show that intervention combinations can exhibit positive complementarity ($\gamma > 0$), independence ($\gamma = 0$), or negative complementarity ($\gamma < 0$), with important implications for portfolio design. The framework provides testable predictions and has been validated in field applications across retirement savings, energy conservation, and health behavior domains.

\vspace{1em}
\noindent\textit{JEL Classification:} D91, D90, C93, H31\\
\textit{Keywords:} Behavioral interventions, nudges, choice architecture, complementarity, present bias
\end{abstract}

\newpage

% =============================================================================
% 1. INTRODUCTION
% =============================================================================
\section{Introduction}

The application of behavioral economics to policy and organizational design has grown substantially over the past two decades. Following the seminal work of Thaler and Sunstein (2008) on ``nudges'' and the establishment of behavioral insights units in governments worldwide, behavioral interventions have become a standard tool in the policy maker's toolkit. Yet despite billions invested and thousands of interventions deployed, systematic evidence suggests that the majority of behavioral interventions underperform expectations (DellaVigna and Linos, 2022).

We argue that this systematic underperformance is not due to fundamental limitations of behavioral science, but rather to the absence of rigorous methodology for intervention design. Current practice is characterized by ad-hoc selection of intervention types, inadequate consideration of contextual factors, and failure to account for interaction effects when interventions are combined. The result is a field that, despite strong theoretical foundations, struggles to achieve consistent results in application.

This paper presents a systematic framework for behavioral intervention design that addresses these limitations. Our approach builds on three decades of experimental research in behavioral economics, particularly the work on social preferences (Fehr and Schmidt, 1999; Fehr and Fischbacher, 2002), present bias (Laibson, 1997; O'Donoghue and Rabin, 1999), and choice architecture (Madrian and Shea, 2001; Thaler and Benartzi, 2004).

The framework makes three primary contributions. First, we develop a comprehensive taxonomy of nine intervention types, each characterized by a distinct behavioral mechanism through which it affects decision-making. This taxonomy provides a common language for intervention design and enables systematic comparison across applications.

Second, we introduce a 20-field schema for intervention specification that captures all dimensions relevant to effectiveness prediction. This schema serves as a standardized documentation format that enables organizational learning---a critical gap in current practice where knowledge from past interventions is rarely preserved or transferred.

Third, we establish five axioms governing effective intervention design. These axioms formalize the conditions under which interventions succeed or fail, providing testable predictions for applied work.

The remainder of this paper is organized as follows. Section 2 presents the theoretical foundations, including the nine intervention types and their mechanisms. Section 3 develops the formal model of intervention effectiveness. Section 4 presents the five design axioms with supporting evidence. Section 5 discusses applications and extensions. Section 6 concludes.


% =============================================================================
% 2. THEORETICAL FOUNDATIONS
% =============================================================================
\section{Theoretical Foundations}

\subsection{The Intervention Design Problem}

Consider a principal (government, organization, or individual) seeking to change the behavior of a target population. The principal can deploy one or more interventions from a set of available options. The fundamental design problem is to select and configure interventions to maximize behavioral change subject to resource constraints.

Formally, let $\mathcal{I} = \{I_1, I_2, \ldots, I_n\}$ denote the set of available interventions. Each intervention $I_i$ is characterized by a vector of attributes including its type $\tau_i$, target context $\psi_i$, and implementation parameters $\theta_i$. The effectiveness of intervention $I_i$ in context $\Psi$ for population segment $s$ is given by:

\begin{equation}
E(I_i | \Psi, s) = E_{\text{base}}(I_i) \cdot \phi(\Psi) \cdot \sigma_s
\label{eq:effectiveness}
\end{equation}

\noindent where $E_{\text{base}}(I_i)$ represents the baseline effectiveness under ideal conditions, $\phi(\Psi)$ is a context adjustment factor, and $\sigma_s$ is a segment-specific multiplier.

This formulation highlights the three key dependencies that current practice typically neglects: context sensitivity, segment heterogeneity, and (as we develop below) complementarity effects when interventions are combined.

\subsection{A Taxonomy of Nine Intervention Types}

We identify nine fundamental intervention types, distinguished by their underlying behavioral mechanisms. This taxonomy is exhaustive in the sense that any behavioral intervention can be classified into one of these categories or as a hybrid combining multiple types.

\begin{definition}[Intervention Type]
An intervention type $\tau$ is defined by its primary behavioral mechanism---the psychological or structural pathway through which it affects decision-making.
\end{definition}

\textbf{Type 1: Informative Interventions.} These interventions operate through knowledge provision, addressing information asymmetries that prevent optimal decision-making. Examples include nutrition labels, financial literacy programs, and health risk communications. The underlying assumption is that the target population would change behavior if properly informed.

\textbf{Type 2: Normative Interventions.} These leverage social preferences by communicating descriptive norms (what others do) or injunctive norms (what others approve of). The extensive experimental literature on social preferences (Fehr and Fischbacher, 2002) establishes that humans are not purely self-interested; rather, they care about others' behavior and approval. Normative interventions activate these preferences.

\textbf{Type 3: Structural Interventions.} These modify the choice environment to make desired behaviors easier and undesired behaviors harder. The canonical example is default enrollment in retirement savings (Madrian and Shea, 2001), which exploits the fact that most individuals are ``passive'' with respect to default options (Chetty et al., 2014).

\textbf{Type 4: Incentive-Based Interventions.} These directly modify the financial costs and benefits of behavior through taxes, subsidies, or performance payments. While straightforward in principle, incentive interventions interact in complex ways with intrinsic motivation, potentially ``crowding out'' pro-social behavior (Frey and Jegen, 2001).

\textbf{Type 5: Commitment Interventions.} These address present bias ($\beta < 1$) by enabling individuals to constrain their future choices. The effectiveness of commitment devices depends critically on whether the individual is a ``sophisticate'' who recognizes their self-control problem (O'Donoghue and Rabin, 1999).

\textbf{Type 6: Feedback Interventions.} These provide real-time information about behavior and its consequences. The Opower energy reports studied by Allcott (2011) represent a canonical example, combining feedback with social comparison to achieve persistent energy reductions.

\textbf{Type 7: Identity-Based Interventions.} These activate or reshape self-concept to align behavior with identity. Bryan et al. (2011) demonstrate that framing voting as being ``a voter'' (identity) rather than ``voting'' (action) significantly increases turnout.

\textbf{Type 8: Temporal Interventions.} These use time strategically through deadlines, cooling-off periods, and optimal timing of prompts. The interaction between temporal interventions and present bias creates complex dynamics that require careful analysis.

\textbf{Type 9: Hybrid Interventions.} Many real-world interventions combine multiple mechanisms. The analysis of hybrid interventions requires understanding complementarity effects, which we develop in Section 3.

Table \ref{tab:types} summarizes the nine types with their primary mechanisms and optimal application contexts.

\begin{table}[H]
\centering
\caption{Taxonomy of Intervention Types}
\label{tab:types}
\small
\renewcommand{\arraystretch}{1.2}
\begin{tabular}{@{}clll@{}}
\toprule
\textbf{Type} & \textbf{Name} & \textbf{Mechanism} & \textbf{Optimal Context} \\
\midrule
1 & Informative & Knowledge provision & Information gaps \\
2 & Normative & Social preferences & High social orientation \\
3 & Structural & Choice architecture & Passive populations \\
4 & Incentive & Financial modification & Price-sensitive decisions \\
5 & Commitment & Pre-commitment & Present-biased sophisticates \\
6 & Feedback & Real-time information & Behavior monitoring \\
7 & Identity & Self-concept activation & Malleable identity \\
8 & Temporal & Strategic timing & Time-sensitive contexts \\
9 & Hybrid & Multiple mechanisms & Complex behaviors \\
\bottomrule
\end{tabular}
\end{table}


% =============================================================================
% 3. THE FORMAL MODEL
% =============================================================================
\section{The Formal Model}

\subsection{Intervention Specification}

We propose a standardized schema for intervention specification consisting of 20 fields organized into three tiers of increasing detail.

\begin{definition}[Light Mode Specification]
The minimum viable specification consists of six fields: (F1) Title, (F2) Type $\tau \in \{1, \ldots, 9\}$, (F3) Context level $\psi$, (F4) Target utility structure, (F5) Welfare dimension, and (F6) Journey phase $\phi$.
\end{definition}

The Light Mode specification captures the essential information required to evaluate an intervention's appropriateness for a given application. We demonstrate in Section 4 that most intervention failures can be attributed to violations involving these core fields.

The Hybrid Mode (F1--F12) adds framing logic, autonomy level, target segment, complementarity links, side effect risk, and time horizon. The Profound Mode (F1--F18) provides research-grade documentation including evaluation plans, path functions, and system integration parameters.

\subsection{The Complementarity Structure}

A critical insight of this framework is that intervention effects are not additive. When interventions are combined, their joint effect may exceed (positive complementarity), equal (independence), or fall short of (negative complementarity) the sum of individual effects.

\begin{definition}[Complementarity Coefficient]
For interventions $I_i$ and $I_j$, the complementarity coefficient $\gamma_{ij} \in [-1, 1]$ measures the deviation from additivity in combined effects.
\end{definition}

\begin{proposition}[Combined Effectiveness]
The effectiveness of combining interventions $I_i$ and $I_j$ is given by:
\begin{equation}
E(I_i \cup I_j) = E(I_i) + E(I_j) + \gamma_{ij} \cdot \sqrt{E(I_i) \cdot E(I_j)}
\label{eq:complementarity}
\end{equation}
where $\gamma_{ij} > 0$ indicates synergy, $\gamma_{ij} = 0$ indicates independence, and $\gamma_{ij} < 0$ indicates interference.
\end{proposition}

The complementarity coefficient has important empirical content. For example, the extensive literature on motivation crowding (Frey and Jegen, 2001; Gneezy and Rustichini, 2000) suggests that financial incentives can undermine intrinsic motivation, implying $\gamma < 0$ for combinations of incentive-based and normative interventions.

Table \ref{tab:complementarity} presents estimated complementarity coefficients for common intervention pairs based on meta-analytic evidence.

\begin{table}[H]
\centering
\caption{Estimated Complementarity Coefficients}
\label{tab:complementarity}
\small
\renewcommand{\arraystretch}{1.2}
\begin{tabular}{@{}lcc@{}}
\toprule
\textbf{Intervention Pair} & \textbf{$\gamma_{ij}$} & \textbf{Source} \\
\midrule
Default + Escalation & $+0.6$ & Thaler and Benartzi (2004) \\
Information + Social norms & $+0.4$ & Allcott (2011) \\
Feedback + Commitment & $+0.4$ & Ashraf et al. (2006) \\
Social norms + Feedback & $+0.5$ & Allcott and Rogers (2014) \\
Incentive + Social norms & $-0.3$ & Gneezy and Rustichini (2000) \\
Commitment + Incentive & $-0.2$ & Charness and Gneezy (2009) \\
\bottomrule
\end{tabular}
\end{table}

\subsection{The Journey Phase Model}

Behavioral change follows a predictable sequence of phases. We adopt a five-phase model: Awareness, Trigger, Action, Maintenance, and Stabilization. Different intervention types are optimally suited to different phases.

\begin{proposition}[Journey-Type Alignment]
Let $\phi \in \{A, T, C, M, S\}$ denote the journey phase and $\tau$ the intervention type. The effectiveness multiplier $\mu(\tau, \phi)$ satisfies:
\begin{equation}
\mu(\tau, \phi) = \begin{cases}
1.0 & \text{if } \tau \text{ is aligned with } \phi \\
0.3 & \text{if } \tau \text{ is misaligned with } \phi
\end{cases}
\end{equation}
\end{proposition}

This proposition formalizes the intuition that intervention timing matters. An informative intervention applied in the Action phase (when awareness already exists) achieves at most 30\% of its potential effectiveness.

\subsection{Segment Heterogeneity}

The same intervention can have dramatically different effects on different population segments. Let $s \in \mathcal{S}$ denote a behavioral segment characterized by preference parameters (social orientation, present bias, autonomy needs, etc.).

\begin{proposition}[Segment Multiplier]
The segment-specific effectiveness is:
\begin{equation}
E(I | s) = E_{\text{base}}(I) \cdot \sigma_s
\end{equation}
where $\sigma_s \in [-0.5, 2.0]$. Negative values indicate reactance effects.
\end{proposition}

The possibility of negative segment multipliers is empirically documented. Normative interventions can trigger reactance in autonomy-oriented individuals, leading to behavior change in the opposite of the intended direction (Schultz et al., 2007).


% =============================================================================
% 4. THE FIVE DESIGN AXIOMS
% =============================================================================
\section{The Five Design Axioms}

We now present five axioms that govern effective intervention design. These axioms synthesize the theoretical framework and provide actionable guidance for practitioners.

\begin{axiom}[Completeness]
An intervention can only be evaluated or deployed if all required fields for its application depth are specified.
\end{axiom}

This axiom establishes a minimum documentation standard. Interventions with missing specifications cannot be properly evaluated, and their implementation outcomes cannot contribute to organizational learning.

\begin{axiom}[Context Dependency]
No intervention has context-independent effectiveness. The same intervention achieves different results in different contexts.
\end{axiom}

The extensive literature on external validity in behavioral economics supports this axiom (Allcott, 2015). Effect sizes observed in laboratory settings or specific field contexts do not automatically transfer to new settings.

\begin{axiom}[Journey Phase Alignment]
Intervention effectiveness is maximized when the intervention type matches the target population's actual journey phase. Misaligned interventions achieve less than 30\% of potential effectiveness.
\end{axiom}

This axiom has direct operational implications. Before selecting an intervention type, practitioners must diagnose the journey phase of the target population. Commitment devices fail in the Awareness phase; informative interventions are redundant in the Action phase.

\begin{axiom}[Complementarity Structure]
Intervention pairs can reinforce each other ($\gamma > 0$), have independent effects ($\gamma = 0$), or undermine each other ($\gamma < 0$). Portfolio design must account for complementarity.
\end{axiom}

The complementarity axiom explains why intervention portfolios sometimes underperform expectations even when individual interventions have proven effectiveness. Negative complementarity can reduce combined effects below what either intervention would achieve alone.

\begin{axiom}[Segment Specificity]
The same intervention can have dramatically different effects on different behavioral segments, including opposite effects.
\end{axiom}

This axiom implies that average treatment effects can be misleading. A normative intervention with zero average effect might have strong positive effects on socially-oriented segments and strong negative effects (reactance) on autonomy-oriented segments.


% =============================================================================
% 5. APPLICATIONS AND EVIDENCE
% =============================================================================
\section{Applications and Evidence}

\subsection{Retirement Savings}

The retirement savings domain provides the most extensive evidence for systematic intervention design. Madrian and Shea (2001) demonstrated that automatic enrollment (Type 3: Structural) increased participation from 49\% to 86\%---a 37 percentage point improvement with minimal implementation cost.

Thaler and Benartzi (2004) extended this with the Save More Tomorrow program, combining default enrollment with automatic escalation. The complementarity coefficient of $\gamma \approx 0.6$ between these interventions (Table \ref{tab:complementarity}) explains why the combined program achieved savings rates that neither intervention could achieve alone.

Chetty et al. (2014) provide crucial evidence on segment heterogeneity. They find that 85\% of savers are ``passive''---accepting default options without active decision-making. For this majority, structural interventions are highly effective ($\sigma_s \approx 1.5$). The remaining 15\% are ``active savers'' who optimize regardless of defaults ($\sigma_s \approx 0$).

\subsection{Energy Conservation}

Allcott (2011) and Allcott and Rogers (2014) provide rigorous evidence on normative interventions (Type 2) in energy conservation. The Opower home energy reports, which combine feedback (Type 6) with social comparison (Type 2), achieved persistent 2--3\% energy reductions.

The positive complementarity between feedback and social norms ($\gamma \approx 0.5$) is evident in the program's success. Social norms provide motivation; feedback enables action. Neither alone achieves the combined effect.

However, the same studies document segment heterogeneity. High-use households show larger treatment effects ($\sigma_s \approx 1.3$), while low-use households show minimal response ($\sigma_s \approx 0.4$). Importantly, the ``boomerang effect''---where low-use households increase consumption upon learning they use less than neighbors---demonstrates the possibility of negative segment effects.

\subsection{Health Behaviors}

The commitment device literature in health contexts illustrates Axiom 3 (Journey Phase Alignment). Commitment contracts for smoking cessation (Gine et al., 2010) and weight loss (Volpp et al., 2008) show effectiveness only when participants are already motivated to change---that is, in the Trigger or Action phase.

Attempting commitment interventions in the Awareness phase fails because individuals who do not yet recognize their behavior as problematic will not voluntarily constrain their future choices. This explains the consistently low uptake of commitment devices in general populations.


% =============================================================================
% 6. CONCLUSION
% =============================================================================
\section{Conclusion}

This paper presents a systematic framework for behavioral intervention design. Our framework addresses the fundamental question of why behavioral interventions so often underperform expectations: it is not that behavioral science is wrong, but that intervention design practice lacks the rigor that behavioral science itself demands.

The framework makes three contributions. First, the nine-type taxonomy provides a common language for intervention classification. Second, the 20-field schema enables standardized documentation that supports organizational learning. Third, the five axioms establish testable conditions for intervention success.

Several directions merit further research. First, the complementarity coefficients in Table \ref{tab:complementarity} represent current best estimates; more precise estimation requires systematic meta-analysis across applications. Second, the journey phase model requires validation in additional behavioral domains. Third, the segment multipliers require development of practical diagnostic tools for field applications.

Despite these limitations, the framework provides actionable guidance for practitioners. The core message is simple: systematic design matters. Before deploying any behavioral intervention, practitioners should (i) classify the intervention type, (ii) verify journey phase alignment, (iii) assess segment composition, and (iv) check complementarity with existing interventions. This discipline will not guarantee success, but it will substantially improve the odds.

\vspace{2em}

% =============================================================================
% REFERENCES
% =============================================================================
\section*{References}
\small

\begin{description}[style=nextline, leftmargin=0pt, itemsep=0.5em]

\item[Allcott, Hunt. 2011.] ``Social Norms and Energy Conservation.'' \textit{Journal of Public Economics} 95(9--10): 1082--1095.

\item[Allcott, Hunt. 2015.] ``Site Selection Bias in Program Evaluation.'' \textit{Quarterly Journal of Economics} 130(3): 1117--1165.

\item[Allcott, Hunt, and Todd Rogers. 2014.] ``The Short-Run and Long-Run Effects of Behavioral Interventions.'' \textit{American Economic Review} 104(10): 3003--3037.

\item[Ashraf, Nava, Dean Karlan, and Wesley Yin. 2006.] ``Tying Odysseus to the Mast.'' \textit{Quarterly Journal of Economics} 121(2): 635--672.

\item[Bryan, Christopher J., Gregory M. Walton, Todd Rogers, and Carol S. Dweck. 2011.] ``Motivating Voter Turnout by Invoking the Self.'' \textit{Proceedings of the National Academy of Sciences} 108(31): 12653--12656.

\item[Charness, Gary, and Uri Gneezy. 2009.] ``Incentives to Exercise.'' \textit{Econometrica} 77(3): 909--931.

\item[Chetty, Raj, John N. Friedman, S{\o}ren Leth-Petersen, Torben Heien Nielsen, and Tore Olsen. 2014.] ``Active vs. Passive Decisions and Crowd-Out in Retirement Savings Accounts.'' \textit{Quarterly Journal of Economics} 129(3): 1141--1219.

\item[DellaVigna, Stefano, and Elizabeth Linos. 2022.] ``RCTs to Scale: Comprehensive Evidence From Two Nudge Units.'' \textit{Econometrica} 90(1): 81--116.

\item[Fehr, Ernst, and Urs Fischbacher. 2002.] ``Why Social Preferences Matter---The Impact of Non-Selfish Motives on Competition, Cooperation and Incentives.'' \textit{Economic Journal} 112(478): C1--C33.

\item[Fehr, Ernst, and Klaus M. Schmidt. 1999.] ``A Theory of Fairness, Competition, and Cooperation.'' \textit{Quarterly Journal of Economics} 114(3): 817--868.

\item[Frey, Bruno S., and Reto Jegen. 2001.] ``Motivation Crowding Theory.'' \textit{Journal of Economic Surveys} 15(5): 589--611.

\item[Gine, Xavier, Dean Karlan, and Jonathan Zinman. 2010.] ``Put Your Money Where Your Butt Is.'' \textit{American Economic Journal: Applied Economics} 2(4): 213--235.

\item[Gneezy, Uri, and Aldo Rustichini. 2000.] ``A Fine Is a Price.'' \textit{Journal of Legal Studies} 29(1): 1--17.

\item[Laibson, David. 1997.] ``Golden Eggs and Hyperbolic Discounting.'' \textit{Quarterly Journal of Economics} 112(2): 443--478.

\item[Madrian, Brigitte C., and Dennis F. Shea. 2001.] ``The Power of Suggestion.'' \textit{Quarterly Journal of Economics} 116(4): 1149--1187.

\item[O'Donoghue, Ted, and Matthew Rabin. 1999.] ``Doing It Now or Later.'' \textit{American Economic Review} 89(1): 103--124.

\item[Schultz, P. Wesley, Jessica M. Nolan, Robert B. Cialdini, Noah J. Goldstein, and Vladas Griskevicius. 2007.] ``The Constructive, Destructive, and Reconstructive Power of Social Norms.'' \textit{Psychological Science} 18(5): 429--434.

\item[Thaler, Richard H., and Shlomo Benartzi. 2004.] ``Save More Tomorrow.'' \textit{Journal of Political Economy} 112(S1): S164--S187.

\item[Thaler, Richard H., and Cass R. Sunstein. 2008.] \textit{Nudge: Improving Decisions About Health, Wealth, and Happiness}. Yale University Press.

\item[Volpp, Kevin G., et al. 2008.] ``Financial Incentive--Based Approaches for Weight Loss.'' \textit{JAMA} 300(22): 2631--2637.

\end{description}

\end{document}
